%==============================================================================
%  CHAPITRE 4 - Pipeline DevOps (CI/CD)
%==============================================================================
\chapter{Pipeline DevOps (CI/CD)}
\label{chap:devops}

La chaîne CI/CD s'appuie sur \textbf{GitHub Actions} (ou GitLab CI au choix)
et se déclenche à chaque \emph{push} ou \emph{pull request}.

\section{Environnements}

\begin{center}
\small
\begin{tabular}{llll}
\toprule
\textbf{Environnement} & \textbf{Branche} & \textbf{URL} & \textbf{Usage} \\
\midrule
Local      & \texttt{feature/*}  & \texttt{http://localhost:8080/zumm} & Développement local (Spring Boot) \\
Dev        & \texttt{develop}    & \texttt{https://dev.zumm.local}     & Intégration continue, tests \\
Staging    & \texttt{release/*}  & \texttt{https://staging.zumm.local} & Recette, démo sprint review \\
Production & \texttt{main}       & \texttt{https://zumm.local}         & Déploiement manuel validé \\
\bottomrule
\end{tabular}
\end{center}

\section{Vue d'ensemble du pipeline}

\begin{figure}[ht]
\centering
\begin{tikzpicture}[
  node distance=4mm and 5mm,
  stage/.style={draw=mielfonce, fill=grisdoux, rounded corners=2pt,
    minimum width=2.45cm, minimum height=1.15cm, align=center,
    font=\scriptsize\bfseries, text=bleunuit},
  manual/.style={stage, draw=vertsauge, dashed},
  fleche/.style={-{Stealth[length=2.4mm]}, thick, color=miel}
]
\node[stage] (build)   {BUILD\\ \scriptsize Maven (3 min)};
\node[stage, right=of build] (test)    {TESTS\\ \scriptsize JUnit + Testcontainers\\ \scriptsize (13 min)};
\node[stage, right=of test] (qualite) {QUALITÉ\\ \scriptsize Sonar + OWASP\\ \scriptsize (9 min)};
\node[stage, right=of qualite] (docker)  {DOCKER\\ \scriptsize Image (6 min)};
\node[stage, below=of docker] (staging) {STAGING\\ \scriptsize Auto (2 min)};
\node[stage, left=of staging] (e2e)     {E2E + k6\\ \scriptsize (25 min)};
\node[manual, left=of e2e] (prod)    {PRODUCTION\\ \scriptsize Manuel\\ \scriptsize (approval gate)};
\draw[fleche] (build) -- (test);
\draw[fleche] (test) -- (qualite);
\draw[fleche] (qualite) -- (docker);
\draw[fleche] (docker) -- (staging);
\draw[fleche] (staging) -- (e2e);
\draw[fleche] (e2e) -- (prod);
\end{tikzpicture}
\caption{Enchaînement des étapes du pipeline CI/CD (durée totale $\approx$ 1 h).}
\label{fig:pipeline}
\end{figure}

\section{Détail des étapes}

\begin{center}
\small
\setlength{\tabcolsep}{4pt}%
\begin{tabular}{p{2.3cm}p{4.4cm}p{4.7cm}p{3.0cm}c}
\toprule
\textbf{Étape} & \textbf{Description} & \textbf{Commandes} & \textbf{Seuils / conditions} & \textbf{Durée} \\
\midrule
Build & Compilation Maven + packaging JAR autoporteur (JDK 17, Maven 3.9+, Spring Boot 3)
  & \texttt{mvn clean compile}\newline \texttt{mvn package -DskipTests}
  & --- & 3 min \\
\addlinespace
Tests unitaires & JUnit 5 + Mockito + JaCoCo
  & \texttt{mvn test}\newline \texttt{mvn jacoco:report}
  & couverture $\geq 70\,\%$, 0 échec & 5 min \\
\addlinespace
Tests d'intégration & Testcontainers : PostgreSQL réel (PostGIS + TimescaleDB)
  & \texttt{mvn verify -P integration-tests}
  & --- & 8 min \\
\addlinespace
Analyse qualité & SonarQube + SpotBugs + Checkstyle
  & \texttt{mvn sonar:sonar}\newline \texttt{mvn spotbugs:check}\newline \texttt{mvn checkstyle:check}
  & 0 bug, 0 vulnérabilité, $<50$ \emph{code smells} & 5 min \\
\addlinespace
Sécurité & Scan des dépendances OWASP
  & \texttt{mvn dependency-check:check}
  & --- & 4 min \\
\addlinespace
Build Docker & Image multi-étapes (build Maven puis JRE + JAR Zümm)
  & \texttt{docker build -t zumm:\$\{VERSION\} .}
  & --- & 6 min \\
\addlinespace
Déploiement staging & Déploiement automatique
  & \texttt{docker-compose -f docker-compose.staging.yml up -d}
  & branches \texttt{release/*}, \texttt{main} & 2 min \\
\addlinespace
Tests E2E & Selenium / Cypress sur staging
  & \texttt{mvn test -P e2e}
  & --- & 10 min \\
\addlinespace
Tests de charge & k6, performance et robustesse
  & \texttt{k6 run tests/load/}
  & p95 $<$ 500 ms, erreurs $<1\,\%$ & 15 min \\
\addlinespace
Déploiement production & Déploiement \textbf{manuel} (\emph{approval gate})
  & \texttt{docker-compose -f docker-compose.prod.yml up -d}
  & branche \texttt{main} + approbation & 2 min \\
\bottomrule
\end{tabular}
\end{center}

\section{Infrastructure as Code}

\begin{itemize}
  \item \textbf{Docker} : \texttt{Dockerfile} applicatif multi-étapes (JAR Spring Boot),
        \texttt{Dockerfile} base de données (PostgreSQL + PostGIS),
        \texttt{docker-compose.yml} pour la stack complète
        (application + PostgreSQL + Nginx + monitoring) ;
  \item \textbf{Kubernetes} (optionnel, scaling futur) : manifestes
        \texttt{deployment.yml}, \texttt{service.yml}, \texttt{ingress.yml},
        \texttt{configmap.yml}.
\end{itemize}

\section{Mise en place progressive}

Le pipeline complet décrit ci-dessus est la cible ; sa mise en place suit
la montée en puissance du projet :

\begin{itemize}
  \item \textbf{CI dès le sprint 1} (build + tests + Sonar), mais CD staging
        seulement à partir du sprint 2 : inutile de maintenir quatre
        environnements avant d'avoir un WAR stable ;
  \item \textbf{Quality gate progressif} : exiger $70\,\%$ de couverture dès
        le sprint 1 est irréaliste sur du code d'infrastructure --- commencer
        à $50\,\%$ et monter de 5 points par sprint jusqu'à la cible ;
  \item \textbf{Gestion des secrets} (GitHub Secrets / vault) intégrée au
        pipeline avant le sprint 4 : l'authentification OAuth (US-020) en
        dépend ;
  \item \textbf{Tests E2E et de charge} activés à partir du sprint 5, quand
        les parcours utilisateur sont stabilisés.
\end{itemize}

\section{Stratégie de branches Git}

\begin{center}
\begin{tabular}{ll}
\toprule
\textbf{Branche} & \textbf{Rôle} \\
\midrule
\texttt{main}       & Production stable \\
\texttt{develop}    & Intégration continue \\
\texttt{release/*}  & Stabilisation d'une release \\
\texttt{feature/*}  & Développement d'une fonctionnalité \\
\texttt{hotfix/*}   & Correction urgente en production \\
\bottomrule
\end{tabular}
\end{center}

\begin{encadre}
\textbf{Fichiers opérationnels.} Les fichiers exécutables du pipeline
(\texttt{github-actions.yml}, \texttt{Dockerfile},
\texttt{docker-compose.yml}) sont maintenus dans
\path{roadmap/operationnel/03_devops_pipeline/} ;
ce chapitre en
décrit le contrat (étapes, seuils, conditions).
\end{encadre}
